\section{Open Problem}

One possible question is that if $\sigma$ is bounded, then in term of the solution $u_{\epsilon}$ of the stochastic pde
\[
\partial_t u_{\epsilon} = \tfrac{1}{2}\Delta u_{\epsilon} + \sigma(u)dW^{\epsilon},
\]
what will $u_{\epsilon}$ converges to as $\epsilon \to 0$.

\subsection{Limit behavior of the parabolic equation}

In this part, we consider a quasi-linear parabolic equation
\[
u_t = \tfrac12 uu_{xx},\quad t\in [0,\infty),x\in\R
\]
and according to the paper by Alex.D and Cole.G, we have such conclusions about the existence and regularity of the solution of this equation. Firstly, let us define the weak solution for the equation

\begin{definition}[Weak solution]
Given $T>0$, we say that
\(
u \in \mathcal{C}\big([0,T);W^{1,\infty}_{\mathrm{loc}}(\R)\big)
\)
is a \emph{weak solution} of the parabolic equation on $[0,T)$ if $u\ge 0$ and for all
$\psi\in C_c^{\infty}([0,T)\times \R)$,
\begin{equation}
\int_{\R} u_0\psi(0,\cdot)
+ \int_{[0,T)\times\R}
\left[u\,\partial_t\psi - \frac12 \,\partial_x(\psi u)\,\partial_x u\right] = 0.
\label{eq:weak}
\end{equation}
\end{definition}

With better regularity of $u$ and $\partial_x^2 u$, we have

\begin{definition}[Strong solution]
Given $\overline T>0$, a weak solution of the parabolic equation on $[0,\overline T)$ is called
\emph{strong} if for every $T\in (0,\overline T)$ we have
\[
u \in \langle x\rangle^{2}L^{\infty}([0,T]\times\R)
\qquad\text{and}\qquad
\partial_x^2 u \in L^1\big([0,T];L^{\infty}(\R)\big).
\]
where
\[
\langle x\rangle = \sqrt{(x^2+1)}
\]
is the Japanese braacket.
\end{definition}

For merely weak solution, we have a lemma to secure the smoothness

\begin{lemma}
    Let $u$ be a weak solution of the parabolic equation on a time interval $[0,  T)$. Then on the set $\{(t, x) \in (0, T) \times \R | u(t,x) > 0\}$, $u$  is qualitatively smooth and satisfies the parabolic equation pointwise.
\end{lemma}

By the way, the equation admits some kind of comparsion principle

\begin{proposition}
    Given $T > 0$, suppose $u^{\pm}\in \mathcal{C}_{\text{loc}}([0,T]\times \R)$ satisfy $\partial_x^2 u^{\pm} \in L^1_tL^{\infty}_x$ and $\partial_t u \in \langle x\rangle^2L^1_tL^{\infty}_x$ as well as
    \[
    \partial_t u^+ \geq \tfrac12 u^+\partial_x^2 u^+,\quad \partial_t u^- \leq \tfrac12 u^-\partial_x^2 u^-, \quad a.e.
    \]
    Then if $u^-(0,\cdot) \leq u^+(0,\cdot)$, we have $u^- \leq u^+$.
\end{proposition}

\subsubsection{Main Conclusion}

Now we would like to consider for some initial condition with relatively good properties, if the limit of $u$ at the time of infinity will exist and what it is, and a basic observation is that $u(t,\cdot)$ will approach to some constant function on $\R$ in general, so the we have a claim

\begin{theorem}[Limit Behavior]
    Assume $u \in \mathcal{C}_{\text{loc}}([0,\infty)\times\R)$ is a weak solution of the parabolic equation, with a positive bounded initial condition $u_0 = u(0,\cdot)$ and $\partial^2_x u \in L_t^1L_x^{\infty}$, and $u_0$ is periodic or
    \[u_0 = f + c\]
    where $c>0$ is a constant and $f$ is compactly supported on $\R$. Then \[\lim_{t\to \infty}u(t,x) = C(u_0)\quad\text{for a.e. }x\in\R,\] where $C(u_0)$ is some constant only related to the initial condition.  
\end{theorem}
\begin{proof}
Step 1: We show that there is some $m,M > 0$ such that $m\leq u \leq M$ for a.e. $(t,x)\in[0,\infty)\times\R$ at first. For both cases of the initial condition, it suffices to show that there exists some $m,M > 0$ and $m \leq u_0 \leq M$ on $[-L,L]$ where $L$ is the period of $u_0$ when it is periodic or $L$ is large enough such that $\text{supp}\ f\subset [-L,L]$ in the other case. We may know $\inf_{[-L,L]} u_0 > 0$ and by the boundedness of $u_0$ we know there are $m,M>0$ such that $m \leq u_0 \leq M$.\par Now let us apply the strong maximum(minimum) principle to $u$, and lets just consider $s = \sup_t\{\inf u(t,\cdot) \geq m/2\}$ and since $[-L,L]$ is compact, we may know $s>0$. If $s < \infty$, then $u$ is postive on $[0,s)\times[-L,L]$ and by the strong minimum principle, $u$ will obtain its minimum on $\{0\}\times[-L,L]\cup[0,s)\times \{-L,L\}$, however, when $u$ is periodic and the minimum is obtained on $[0,s)\times[-L,L]$, then there exists an minimum in the interior of $[0,s)\times[-L,L]$ by the period and $u$ will become a constant on $[0,s)\times[-L,L]$ which means $u(s,\cdot) \equiv m$ and which is a contradiction. Therefore, the minimum will be obtained on $\{0\}\times [-L,L]$ and hence $u\geq 0$ on $[0,s)\times[-L,L]$. Then we have $u(s,\cdot) \geq m$, which contradicts the definition of $s$ and hence $s = \infty$. An analog argument can be used for $u$ in another case and hence $u$ will have a lower bound $m' > 0$ on the whole $[0,\infty) \times [-L,L]$ and furthermore, on $[0,\infty)\times\R$. Therefore, we may apply the strong maximum and minimum principle for the whole parabolic cylinder and an upper bound can be found similarly.\par
Step 2: Notice here $\partial_t u, \partial_x u,\partial_x^2 u$ are all defined naturally and even smooth by the smoothness of weak solution, so we may consider the energy estimation for $u$. Define
\[
E(t) = \int_{-L}^L (\partial_x u(t,y))^2 dy
\]
and then by compactness of $[-L,L]$ and DCT, we have
\[
\begin{aligned}
    \dfrac{dE(t)}{dt} &= \int_{-L}^L 2\partial_{xt}u(t,y)\partial_x u(t,y) dy\\
    &= 2\left(u_t(t,x)u_x(t,x)|^L_{-L} - \int_{-L}^L u_t(t,y)\partial_{xx}u(t,y)dy\right) \\
    &= - \int^{L}_{-L}u(t,y)\partial_{xx}^2u(t,y)dy \leq 0\\
\end{aligned}
\]
in both cases. So $u(t,\cdot)$ is bounded in $H^1([-L,L])$. Notice $\partial_x^2 u \in L_t^1L^{\infty}$, then by
\[
\lVert\partial_x^2 u\rVert_{L^2([-L,L])} \leq 2L\lVert\partial_x^2 u\rVert_{L^{\infty}(\R)}
\]
and hence
\[
\int_{[0,\infty)}\lVert\partial_x^2 u\rVert_{L^2([-L,L])} \leq 2L \int_{[0,\infty)} \lVert\partial_x^2u\rVert_{L^{\infty}(\R)} < \infty
\]
which means we may choose a subsequence $u_{t_n}$ such that $\lVert\partial_x^2 u(t_n,\cdot)\rVert_{L^{\infty}([-L,L])} \leq 1/n$ and hence $\lVert\partial_x^2 u(t_n,\cdot)\rVert_{L^{2}([-L,L])} \leq 1/n$. By Morrey's inequality, $u(t_n,\cdot),\partial_x(t_n,\cdot)$ are both compactly embedded in $\mathcal{C}^{0,1/2}([-L,L])$ and hence we may choose a further subsequence $u(t_{n_j},\cdot) \to u_{\infty}$ with $\partial_x(t_{n_j},\cdot) \to u'_{\infty}$ in $\mathcal{C}^{0,1/2}([-L,L])$ for some $u_{\infty}$ and $u'_{\infty}$. For any test function $\phi \in C^{\infty}_c([-L,L])$, we have
\[
\begin{aligned}
    \int_{[-L,L]} u_{\infty}\partial_x\phi &= \lim_{j\to\infty} \int_{[-L,L]} u(t_{n_j},\cdot)\partial_x\phi \\ &= -\lim_{j\to\infty} \int_{[-L,L]} \partial_xu(t_{n_j},\cdot)\phi = -\int_{[-L,L]} u'_{\infty}\phi
\end{aligned}
\]
which means $u'_{\infty} = \partial_x u_{\infty}$ in weak sense and hence $u_{\infty} \in W^{1,\infty}([-L,L])$.\par
Step 3: We consider $\partial_x^2 u_{\infty}$ only as a distribution, then for any $\phi \in C^{\infty}_c([-L,L])$, we have
\[
    \begin{aligned}
        \langle \partial_t u (t_n,\cdot), \phi \rangle &= \tfrac12 \int_{[-L,L]}   u (t_n,\cdot)\partial_{x}^2u (t_n,\cdot)\phi \\
        &= \tfrac12 \left(\int_{[-L,L]} \partial_x(u (t_n,\cdot)\partial_x u (t_n,\cdot)) - (\partial_x u (t_n,\cdot))^2 \phi\right) \\
        &= \tfrac14 \int_{[-L,L]} \partial_x^2 (u^2 (t_n,\cdot))\phi - \tfrac12 \int_{[-L,L]}(\partial_x u (t_n,\cdot))^2\phi \\
        &= \tfrac14 \int_{[-L,L]} u^2(t_n,\cdot)\partial_x^2 \phi - \tfrac12 \int_{[-L,L]}(\partial_x u (t_n,\cdot))^2\phi
    \end{aligned}
\]
since $u_{\infty},\partial_xu_{\infty}$ are both continuous function on $[-L,L]$ and $u(t_n,\cdot) \to u_{\infty}, \partial_x u(t_n,\cdot) \to \partial_x u_{\infty}$ uniformly, we have $u^2(t_n,\cdot) \to u^2_{\infty}, (\partial_x u(t_n,\cdot))^2 \to (\partial_x u_{\infty})^2$ uniformly. Therefore we have
\[
\begin{aligned}
     \langle \partial_t u (t_n,\cdot), \phi \rangle &=  \tfrac14 \int_{[-L,L]} u^2(t_n,\cdot)\partial_x^2 \phi - \tfrac12 \int_{[-L,L]}(\partial_x u (t_n,\cdot))^2\phi \\
     &\to \tfrac14 \int_{[-L,L]} u^2_{\infty}\partial_x^2 \phi - \tfrac12 \int_{[-L,L]}(\partial_x u_{\infty})^2\phi \\
     &= \tfrac12 \langle u_{\infty}\partial_x^2 u_{\infty}, \phi\rangle
\end{aligned}
\]
where $u_{\infty}\partial_x^2 u_{\infty}$ is well-defined since $u_{\infty}$ is positive continuous function, which may be approached by smooth function. Notice
\[
\lVert \partial_t u(t_n,\cdot)\rVert_{L^{\infty}(-L,L)} \leq \tfrac1m \lVert\partial_x^2 u(t_n,\cdot)\rVert_{L^{\infty}(-L,L)} \to 0 
\]
and hence $vv_{\infty} = 0$ as a distribution, which means 
\[
\partial_x^2 (u_{\infty}^2) = 2 (\partial_x u_{\infty})^2 \in \mathcal{C}([-L,L])
\]
(Check) We consequently have $\partial_x^2 u_{\infty} \in L^{\infty}([-L,L])$ by some direct computation. Now $u_{\infty}\partial_x^2 u_{\infty}$ may be interpreted by $L^{\infty}([-L,L])$ functions and then $\partial_x^2 u_{\infty} = 0$ by $u_{\infty}$ positive, so $u_{\infty}$ is some constant between $m$ and $M$ when $u_0$ is periodic, and the constant is exactly $c$ as we defined in another case. We denote this constant as $u_{\infty}$ uniformly for the following.
Step 5: Now we may consider the truncation solution $U_T$ which is the solution of the parabolic equation with the initial condition $u(T,\cdot)$, by the comparision principle, we know $U_T$ will be bounded in $[u_{\infty} -\epsilon, u_{\infty} + \epsilon]$ if $|u(T,\cdot) - u_{\infty}|_u < \epsilon$. To sum up, let $T$ to be $t_n$ we chose before and then we will know $u(t+t_n,x) = U_T(t,x) \in [u_{\infty} - \epsilon_n, u_{\infty} + \epsilon_n]$ for any $t\geq 0$ and $\epsilon_n \to 0$ with $n\to\infty$, so we obtained that \[ \lim_{t\to\infty} u(t,\cdot) = u_{\infty}\] in $\mathcal{C}(\R)$, and in particular, for any $x\in \R$, \[\lim_{t\to\infty} u(t,x) = u_{\infty}.\]
\end{proof}